\section{首步分析、击中时间与吸收链}
\label{sec:06-Random-Processes/03-Markov-Chains/05}

你面前摆着一个 Markov 链，你想知道：从状态 \(i\) 出发，最终到达目标集合 \(A\) 的概率是多少？到达之前，平均要等多少步？如果链最终会被几个吸收状态之一捕获，落入每个吸收状态的概率各是多少？

直接模拟你可以跑很多条路径然后统计，但难道没有更干净的办法吗。

你把问题切成两步：先走第一步，到了某个邻居 \(j\)，然后把剩下的问题原封不动地交给同一个未知函数。用方程写下：

\[
\text{未知量}(i)=\text{第一步之后，对邻居的同类未知量取期望}.
\]

这就是\textbf{首步分析}。它不要求任何新工具，只用到条件期望和一次一步的 Markov 性，却一口气覆盖了：

\begin{itemize}
\tightlist
\item
  最终能否到达目标；
\item
  到达前平均需要多少步；
\item
  吸收前访问某状态多少次；
\item
  最终被哪个吸收状态捕获；
\item
  以及带状态相关成本的累计。
\end{itemize}

\subsection{击中概率：一个调和方程}

设目标集合 \(A\)。击中时间定义为

\[
T_A=\inf\{n\ge0:X_n\in A\},
\]

击中概率为

\[
h_i=P_i(T_A<\infty).
\]

如果 \(i\in A\)，答案就是 \(1\)——已经在了，零步击中。如果 \(i\notin A\)，按第一步条件化：

\[
h_i=\sum_jP_{ij}\,h_j.
\]

这里没有什么新技巧，只是把``从 \(i\) 出发最终击中 \(A\) 的概率''分解为``第一步跳到 \(j\) 的概率''乘``从 \(j\) 出发最终击中 \(A\) 的概率''，然后对 \(j\) 求和。未知数 \(h_i\) 在非目标状态上等于邻居值的概率加权平均——这是离散调和函数。边界条件决定具体解。

如果有两个吸收集合 \(A\)（成功）和 \(B\)（失败），想知道``先到 \(A\) 而不是先到 \(B\)''的概率，设置

\[
h_i=1\quad(i\in A),
\qquad
h_i=0\quad(i\in B),
\]

内部仍然满足 \(h=Ph\)。不同的边界把同一个调和方程导向不同的答案。

\subsection{击中时间：每次转移多消耗一步}

现在问不是``会不会''而是``多久''。令

\[
m_i=\E_i[T_A].
\]

在目标内部，\(m_i=0\)——已经到了，不需要再等。在目标外，先走一步消费一个时间单位，然后从新状态继续：

\[
m_i=1+\sum_jP_{ij}\,m_j.
\]

和击中概率方程的唯一区别是多了一个 \(1\)：每次转移，你都在为路径长度这口井里加一滴水。如果每一步停留在状态 \(i\) 还有成本 \(c_i\)（比如等待时间、能耗、手续费），方程自然推广为

\[
v_i=c_i+\sum_jP_{ij}\,v_j.
\]

这些方程只有在击中时间几乎必然有限、且期望有限时才给出所需的唯一有限解。对无限状态链，可能存在多个非负解，或期望发散为无穷。方程写出来容易，解的唯一性和有界性需要额外论证——不要看到线性方程就默认有漂亮答案。

举一个具体场景。客服工单在几个中间状态之间流转：待处理、复核中、已升级。最终要么解决，要么以``未解决''关闭。把解决和未解决关闭设为两个吸收状态，工单从每个中间状态出发最终被解决的概率就是 \(h_i\)，而 \(m_i\) 是结束前的平均流转次数。如果复核和升级耗时不同，还可以用 \(c_i\) 把实际处理时间带入价值方程，算出期望处理总时长。

真正困难的地方往往不是解方程——线性方程组而已——而是确认状态本身是否充分。如果工单的严重程度或已等待时长没有写进状态定义里，下一步流向仍然会依赖更早历史的信息，Markov 递推就不能在这一组状态标签上成立。首步分析的前提不是``写了方程就能解''，而是``状态定义已经包含了所有影响未来的信息''。

\subsection{赌徒破产：把钱输光的概率}

把上面的框架套到一个经典问题。你手上有 \(i\) 元资金，每局以概率 \(p\) 赢一元，以概率 \(q=1-p\) 输一元。资金到达 \(0\) 则破产，到达 \(N\) 则离场。状态空间是 \(0,1,\ldots,N\)，\(0\) 和 \(N\) 是吸收边界。

令 \(h_i\) 为从资金 \(i\) 出发先到 \(N\)（赢到目标）而不是先输光的概率。边界：

\[
h_0=0,\qquad h_N=1.
\]

内部递推：

\[
h_i=ph_{i+1}+qh_{i-1}.
\]

先看公平情形 \(p=q=1/2\)。方程变成 \(h_i=(h_{i+1}+h_{i-1})/2\)，也就是 \(h_i\) 是其左右邻居的算术平均。满足这个条件的必定是线性函数。带入边界 \(h_0=0\)、\(h_N=1\)，直线从 \((0,0)\) 穿过 \((N,1)\)，斜率就是 \(1/N\)：

\[
h_i=\frac iN.
\]

起点资金占目标资金的比例，恰好是率先达到目标的概率。手上有一半的钱，就有一半的胜算。结果简洁到你可以不经计算直接猜出来，但前提是游戏公平。

公平游走的期望吸收时间满足

\[
m_0=m_N=0,
\]
\[
m_i=1+\frac12m_{i+1}+\frac12m_{i-1}.
\]

解为

\[
m_i=i(N-i).
\]

代入验证：\(1+\frac12(i+1)(N-i-1)+\frac12(i-1)(N-i+1)=i(N-i)\)，两边确实相等；边界条件 \(i=0\) 和 \(i=N\) 处也确实为零。

注意这个形状。获胜概率随初始资金线性增长——直观。等待时间却随区间尺度二次增长——不太直观，但可以理解：越靠中间，两个边界都远，链要在两端之间反复游荡很久才肯触碰边界。\(N=100,i=50\) 时，期望等待 2500 步。

当 \(p\neq q\)（游戏不公平），同一递推 \(h_i=ph_{i+1}+qh_{i-1}\) 的通解不再是线性函数。设 \(h_i=r^i\) 代入齐次方程，得特征方程 \(pr^2-r+q=0\)，两个根是 \(1\) 和 \(q/p\)。通解为 \(A+B(q/p)^i\)。结合边界条件 \(h_0=0\)、\(h_N=1\)：

\[
h_i=\frac{1-(q/p)^i}{1-(q/p)^N}.
\]

很小的单步偏差，经过足够多次放大后，会根本性地改变最终命运。如果每局赢的概率只有 \(0.49\)，方程里的 \((q/p)\) 略大于 \(1\)，那个指数让你从任何初始资金出发，先到 \(N\) 的概率都可能低得惊人。

\subsection{吸收链：把递推打包成矩阵}

对有限状态的吸收链，把状态重新排序：暂留状态在前，吸收状态在后。转移矩阵分块为

\[
P=
\begin{bmatrix}
Q&R\\
0&I
\end{bmatrix}.
\]

\(Q\) 是暂留状态之间的转移子矩阵，\(R\) 是从暂留状态一步跳入吸收状态的子矩阵。下半部分 \(0\) 和 \(I\) 表示吸收状态不会再离开。

如果最终以概率一被吸收，基本矩阵

\[
N=(I-Q)^{-1}
=I+Q+Q^2+\cdots
\]

是良定义的。\(N_{ij}\) 是什么？级数展开告诉你：从暂留状态 \(i\) 出发，零步后（如果在 \(j\) 则计一次）、一步后、两步后……落在暂留状态 \(j\) 的概率之和，恰好等于吸收前访问暂留状态 \(j\) 的期望次数。\(I\) 项的存在是因为出发时刻若已在 \(j\) 也算一次访问——是的，你还没迈步就已经``访问''了一次，这取决于你怎么定义计数，实际用的时候以级数和为准。

吸收前期望步数向量：

\[
t=N\mathbf1,
\]

其中 \(\mathbf1\) 是全一向量，右乘 \(N\) 就是把每行的期望访问次数加起来。

进入各吸收状态的概率矩阵：

\[
B=NR.
\]

\(N_{ik}\) 是从 \(i\) 出发访问暂留状态 \(k\) 的期望次数，\(R_{kj}\) 是从 \(k\) 一步跳入吸收状态 \(j\) 的概率，乘起来求和就是从 \(i\) 出发最终被 \(j\) 吸收的概率。这和首步分析的递推方程是同构的——矩阵形式不过是一次性求解了所有起始状态。

\subsection{首步分析的适用范围比矩阵公式大}

\((I-Q)^{-1}\) 这个公式只适用于有限吸收链。首步分析的思维框架能走得更远：

\begin{itemize}
\tightlist
\item
  无限状态随机游走——矩阵无限大，但递推方程的结构照旧；
\item
  状态相关的每一步成本——\(c_i+\sum_jP_{ij}v_j\) 在 \((I-Q)^{-1}\) 公式里没有直接对应，但首步的``当前成本+未来期望''直接写出来；
\item
  到集合前的累积报酬——同样的逻辑加上一个奖励函数；
\item
  连续时间链的生成矩阵方程——加法换成积分，但``当前速率+未来期望''的骨架不变；
\item
  动态规划与 Bellman 方程——首步分析正是 Markov 决策过程价值迭代在零决策情形下的特例。
\end{itemize}

所有这些场景共享同一个骨架：

\[
\text{当前价值} = \text{立即贡献} + \text{下一状态价值的条件期望}.
\]

首步意思是先把第一步的贡献单独拿出来，再将余下无穷多步的期望——恰好等于同一函数的邻居值——交给同一个未知量。方程写下来，边界条件定下来，解就是了。骨架比具体公式更值得记住。
